\chapter{Rationality Block Template and Full Catalogue}
\label{app:rationality}

\section{Purpose}
\label{app:rat:purpose}

Chapter~\ref{ch:security} derives six attack classes, each with a
\emph{Rationality and Scope} block structured per the seven-field
schema of Section~\ref{se:rationality_conventions}. The FC 2027 paper
compresses these blocks into a single matrix
(\texttt{conference/FC-2027/paper/appendix-b-threat-matrix.tex},
§B.2); for the thesis, the block form preserves the full reasoning
chain for each class. A reviewer or deployment engineer can audit
any field against their own environment; a claim that is valid under
the thesis's conditions but not under some downstream reader's
condition should be flagged by a mismatch in fields (3)--(5) rather
than by reading between the lines.

This appendix restates the six blocks in free-text form. The thesis
chapters reference these blocks by subsection label (e.g.,
\S\ref{app:rat:classfee} for class \classfee{}) rather than by
embedding the prose inline.

\section{Class \classfee{}: Fee-Channel Pinning}
\label{app:rat:classfee}

\paragraph{(1) Attacker capability.}
The attacker holds the hot key of the target vault and capital
sufficient for the anchor-descendant chain (approximately
$14{,}300$~sats at the reference dust-relay rate of $3\satvB{}$).
Mempool access is required --- specifically, the ability to broadcast
transactions to a Bitcoin Core node that propagates under default
policy. The attacker is assumed to know the vault's design and axis
values (Kerckhoffs, Section~\ref{ss:kerckhoffs}).

\paragraph{(2) Defender model.}
An online watchtower that monitors the mempool and blockchain and
attempts CPFP replacement on the unvaulting UTXO's anchor output when
an unauthorised trigger is observed. The watchtower operates a
default Bitcoin Core mempool with the $25$-descendant relay limit
active. The attacker's descendant chain saturates this limit before
the watchtower's CPFP transaction can be accepted; BIP-125 Rule 3
then requires the replacement to pay for evicting every descendant,
which is economically infeasible against a dust-rate chain.

\paragraph{(3) Economic preconditions.}
$\mathcal{L}(\mathrm{TM1})$ and the adversary advantage $\mathrm{Adv}$
are both $r$-invariant: the attack cost $C_{\text{pin}} = 25(\rho+1)
\cdot \text{vsize}_{\text{desc}} + \alpha_{\text{anchor}}$ is bounded
in absolute terms (not as a function of vault value or fee rate).
For any reference vault value $V \gg C_{\text{pin}}$, the attack is
rational at every fee regime. We measured $C_{\text{pin}} \approx
14{,}300$~sats against a $0.5$~BTC reference vault; that is $0.03\%$
of the vault.

\paragraph{(4) Protocol / network assumptions.}
Bitcoin Core default mempool policy: $25$-descendant limit from
BIP-125; $3\satvB{}$ dustRelay threshold (policy, not consensus);
RBF Rule 3 requiring the replacement to pay both a higher fee rate
and a higher absolute fee than the cumulative descendants it evicts.
The anchor output sized at $3{,}300$~sats (CTV reference design) is
spendable by any party, which lets the attacker pay its descendant
chain without any key material beyond the hot key used for the
trigger.

\paragraph{(5) Scope of validity.}
Applies iff the covenant's axis \axfee{} has value
\texttt{out-of-band-anchor} --- i.e., fees attach via a CPFP anchor
output rather than being deducted from vault value or supplied via a
separate fee-wallet input. Does not apply to covenants with dynamic
fee inclusion (CCV, OP\_VAULT, CAT+CSFS), whose trigger transactions
do not carry a pinnable anchor handle: the attacker cannot pin a
transaction whose very existence they are trying to delay.

\paragraph{(6) Rational-attacker check.}
Rational at any $r$ given hot-key compromise: payoff $\approx V$
(full vault theft once the watchtower's recovery is blocked and
the hot key completes withdrawal after CSV expiry), cost $\approx
C_{\text{pin}}$, residual value risk to the attacker limited to the
descendant-chain sats.

\paragraph{(7) Counter-factual mitigations.}
Three distinct counter-factuals eliminate or shrink the class:
\begin{enumerate}
\item \textbf{TRUC/BIP-431 (recommended).} The TRUC
  (\texttt{version=3}) transaction type collapses the descendant
  limit to $1$. A \opctv{} unvault that uses TRUC eliminates the
  pinning surface at the relay-policy level; this is available in
  Bitcoin Core v28.0+ and requires only that the unvault template
  commit to the TRUC version.
\item \textbf{Cluster mempool redesign.} The cluster-mempool
  proposal (Wuille and Daftuar~\cite{ClusterMempool}) changes
  replacement-cost accounting in ways that qualitatively alter the
  pinning-cost calculus. Scope-limit rather than a within-scope
  mitigation; treated as future work.
\item \textbf{P2A-style smaller anchor.} Reducing the anchor output
  from $3{,}300$~sats (BIP-118-style bare anchor) to the $330$~sat
  P2A companion specified in BIP-431 reduces the absolute attack
  cost but does not eliminate the structural surface.
\end{enumerate}
A vault-design recommendation derived from this block is:
\emph{always deploy CTV unvault transactions as TRUC wherever Core
v28.0+ is present on the network path}.

\section{Class \classgrief{}: Permissionless Griefing}
\label{app:rat:classgrief}

\paragraph{(1) Attacker capability.}
Zero keys for CCV (the class's only susceptible covenant). The
attacker needs only mempool observation (to detect unauthorised
triggers as they propagate) and the public recovery address.
Kerckhoffs: we assume the recovery address is known, either because
it was published for auditability, reused across vaults under the
same cold key, or revealed by any prior legitimate recovery.

\paragraph{(2) Defender model.}
Any. The attack is pure liveness denial: the attacker does not
extract funds, only forces the owner to re-deposit and re-trigger
repeatedly. The owner's choice is to tolerate the cycle, abandon
the vault, or rotate to a different covenant design. A watchtower
does not help --- the attack triggers the watchtower's own recovery
path, which is what the attacker wants.

\paragraph{(3) Economic preconditions.}
$\mathcal{L}$ scales linearly in fee rate $r$ (each round costs
$\text{vsize}_{\text{def}} \cdot r$ plus the owner's re-trigger
cost). The asymmetry $\gamma = \text{vsize}_{\text{atk}} /
\text{vsize}_{\text{def}}$ is $r$-invariant.

\paragraph{(4) Protocol / network assumptions.}
Mempool-observable trigger transactions (standard under
all Bitcoin Core releases). No RBF dependence: the attacker's
recovery transaction does not contend with the trigger for the same
UTXO spend, only with subsequent withdraw attempts.

\paragraph{(5) Scope of validity.}
Applies iff \axauth{} = \texttt{keyless}. Reduces to a
key-compromise attack (TM6) under any key-gated recovery model;
the key-gated variant has a much narrower attacker pool.

\paragraph{(6) Rational-attacker check.}
Irrational without external incentive: the attacker has no direct
financial gain. Rational only when the external incentive (extortion
payout, competitor disruption, state-actor operation) exceeds the
round-by-round cost.

\paragraph{(7) Counter-factual mitigations.}
\axauth{} = \texttt{recoveryauth-key} (as in OP\_VAULT) eliminates
the class by requiring a Schnorr signature on the recovery leaf;
any key-gated recovery (cold key, hot key, or dedicated authorisation
key) has the same effect. The design tradeoff is that key-gated
recovery reintroduces key-loss risk --- Proposition~1 of
Chapter~\ref{ch:security} formalises the incompatibility.

\section{Class \classscale{}: Per-UTXO Recovery-Cost Scaling}
\label{app:rat:classscale}

\paragraph{(1) Attacker capability.}
The trigger (hot) key of the target vault and capital for $N$
partial-withdrawal trigger transactions. Kerckhoffs: the attacker
knows the vault's axis values, specifically \axrevault{} =
\texttt{yes}.

\paragraph{(2) Defender model.}
A watchtower with recovery authority (keyless for CCV,
\texttt{recoveryauth} for OP\_VAULT). The defender-side degree of
freedom is the batching flag $b \in \{0, 1\}$: single-input
recovery ($b = 0$) scales linearly in UTXO count; batched recovery
($b = 1$) has sub-linear amortisation. Defender policy also matters:
a minimal-disruption policy (reactive per-UTXO recovery) maximises
attacker leverage, while a full-sweep-on-first-alert policy
collapses the class to one round at the cost of operational
disruption.

\paragraph{(3) Economic preconditions.}
Rate-dependent; the single-block-feasibility threshold
$r^\dagger_{\mathrm{TM4}}(V, D, b)$ from
Equation~\ref{eq:r_star} is the governing parameter. Below
$r^\dagger$ the defender absorbs the per-UTXO liability before dust
economics kick in. Above $r^\dagger$ the attacker can drive every
UTXO below the economic abandonment threshold within a single block
window.

\paragraph{(4) Protocol / network assumptions.}
Bitcoin dust thresholds: $546$~sats P2WPKH, $330$~sats P2TR. RBF
Rule 3 does not permit sibling-eviction, so the attacker's
partial-withdrawal chain is single-chain (each partial unvault
consumes the previous output). Minimum relay fee $\approx 1\satvB{}$
bounds the attacker's per-tx cost from below. Mempool min-fee during
congestion can float above the attacker's preferred rate; this is a
scope modifier but not a sufficient mitigation.

\paragraph{(5) Scope of validity.}
Applies iff \axrevault{} = \texttt{yes}. The class is a
\emph{structural cost tail of partial withdrawal}, independent of
adversary behaviour; an attacker accelerates the distribution but
is not required for the class to manifest. Legitimate operational
partial withdrawal can produce the same distribution over time.
Does not apply to CTV or CAT+CSFS (\axrevault{} = \texttt{no}).

\paragraph{(6) Rational-attacker check.}
Rational only above $r^\dagger$ and only with external incentive
(competitor, extortion, state actor). A pure liveness attacker is
irrational at any $r$ --- no financial gain. An attacker who holds
the hot key \emph{and} drives the vault below $r^\dagger$ can extract
value from UTXOs the watchtower abandons, assuming CSV expiry on
those UTXOs.

\paragraph{(7) Counter-factual mitigations.}
Three distinct counter-factuals:
\begin{enumerate}
\item \textbf{Defender batching.} Raises $r^\dagger$ by
  $1.8$--$2.7\times$ (Chapter~\ref{ch:results}). Materially absorbs
  the tail under moderate congestion. Requires BIP-345
  §``Batching'' or BIP-443 default aggregation support in the
  watchtower implementation.
\item \textbf{Full-sweep-on-first-alert.} Collapses the multi-round
  loop to one round at the cost of forcing the owner to re-deposit.
  Appropriate when the operational cost of a full sweep is less than
  the expected loss from the cost tail.
\item \textbf{Operational limits on per-epoch trigger count.} Not a
  protocol-level mitigation, but a watchtower configuration that
  triggers a full sweep after $k$ unauthorised triggers per unit
  time limits the attacker's burst capacity.
\end{enumerate}

\paragraph{Reframing note.}
Earlier framings of this class as ``watchtower exhaustion'' cast it
as a multi-round attacker-vs-watchtower game. That framing does not
survive a defender with full recovery authority who can full-sweep
on first alert, collapsing the multi-round narrative to one round.
The reframe above identifies the class as a structural cost tail of
partial withdrawal --- attacker-accelerable but not
attacker-required --- with the defender's degree of freedom in
recovery batching.

\section{Class \classhot{}: Hot-Key Theft via Destination Substitution}
\label{app:rat:classhot}

\paragraph{(1) Attacker capability.}
The hot (trigger) key; mempool access; optional fee-pinning capital
(if class \classfee{} composes, as it does for CTV).

\paragraph{(2) Defender model.}
An online watchtower able to race the attacker's unvault-to-withdraw
sequence during the CSV delay. The race bias tilts toward the
attacker when class \classfee{} composes (CTV with compromised fee
key); otherwise the defender wins the race generically by
broadcasting the recovery transaction with higher fees.

\paragraph{(3) Economic preconditions.}
Race outcome is $r$-independent except via class \classfee{}
composition, which forces the watchtower's CPFP recovery to fail
regardless of fee rate.

\paragraph{(4) Protocol / network assumptions.}
CSV delay length; mempool propagation dynamics. Riard's replacement
cycling attack~\cite{Ria23} is adjacent but out of scope here; it
would compose with class \classhot{} in the same way that class
\classfee{} composes, but is a mempool-layer rather than a
script-layer mitigation target.

\paragraph{(5) Scope of validity.}
Applies iff \axwdbind{} is one of
\{\texttt{template-hash-digest}, \texttt{per-output-amount+script},
\texttt{dedicated-opcode}\} --- i.e., any value other than
\texttt{signature-dual-bind}. CAT+CSFS, uniquely, has \axwdbind{} =
\texttt{signature-dual-bind} and is immune by construction.

\paragraph{(6) Rational-attacker check.}
Rational only paired with fee-pinning budget (composition with
class \classfee{}) or against a defender with ineffective watchtower
response. Otherwise the defender wins the race generically and the
attack yields griefing only (which is irrational without external
incentive per class \classgrief{}).

\paragraph{(7) Counter-factual mitigations.}
\axwdbind{} = \texttt{signature-dual-bind} eliminates the class by
cryptographic construction. The CAT+CSFS dual-verification pattern
forces the assembled sighash preimage to equal the consensus sighash
preimage under Schnorr EUF-CMA security~\cite{NRS20}: a valid
signature over two distinct messages under the same key would reduce
to discrete-log. The \texttt{hashOutputs} field is a component of
both preimages, so the output set is forced. A hot-key holder
attempting to redirect the withdrawal produces a rejected
transaction, not a theft; the outcome degrades to griefing-only.

\section{Class \classcold{}: Cold-Key Theft via Unconstrained Recovery}
\label{app:rat:classcold}

\paragraph{(1) Attacker capability.}
The cold (recovery) key. Kerckhoffs: we assume compromise, not
inference --- the attack does not derive the cold key from observable
data.

\paragraph{(2) Defender model.}
Any; the attack is instant with no CSV window and no watchtower
intervention window. The cold-key-signed recovery transaction is
valid the moment it is broadcast and there is no script-level check
that would reject an attacker-chosen destination.

\paragraph{(3) Economic preconditions.}
$r$-invariant.

\paragraph{(4) Protocol / network assumptions.}
N/A --- pure signature check; no mempool dynamics. The attack
completes in one transaction with no race window.

\paragraph{(5) Scope of validity.}
Applies iff \axrecbind{} = \texttt{plain-signature}. The only
covenant with this axis value is CAT+CSFS, whose recovery leaf is a
bare \texttt{cold\_pk OP\_CHECKSIG} with \sighash{DEFAULT} and no
script-level output constraint.

\paragraph{(6) Rational-attacker check.}
Rational whenever cold key is compromised; yield $V$ (full vault
theft), cost negligible.

\paragraph{(7) Counter-factual mitigations.}
\axrecbind{} = \texttt{signature-dual-bind} on the recovery leaf
eliminates the class. Poelstra's sketch~\cite{Poe21} describes a
CAT+CSFS-based bound-recovery construction that would move CAT+CSFS
to this axis value; it is not in the reference CAT+CSFS vault
implementation that this thesis measures. Alternative \axrecbind{}
values (\texttt{template-hash-digest} as in CTV,
\texttt{per-output-amount+script} as in CCV,
\texttt{dedicated-opcode} as in OP\_VAULT) also produce immunity.
The class-\classcold{} surface of CAT+CSFS is therefore the
\emph{one structural cost} of its class-\classhot{} immunity, and
it is addressable without moving off CAT+CSFS.

\section{Class \classmode{}: Mode/Parameter Bypass}
\label{app:rat:classmode}

\paragraph{(1) Attacker capability.}
Any observer of a vault whose script contains an undefined mode
byte; no keys required.

\paragraph{(2) Defender model.}
Developer discipline around mode-byte selection; library-provided
named constants (e.g., pymatt's \texttt{CCV\_FLAG\_*}); code review.

\paragraph{(3) Economic preconditions.}
None applicable --- full vault bypass at constant $110\vB{}$/spend
(measured). The class is insensitive to vault value and fee rate.

\paragraph{(4) Protocol / network assumptions.}
BIP-341 and BIP-342 taproot \texttt{OP\_SUCCESSx} semantics: an
undefined opcode value in tapscript terminates execution with
success regardless of stack state. This is the forward-compatibility
mechanism that permits future soft-forks to assign meaning to
currently-undefined bytes without invalidating existing transactions.

\paragraph{(5) Scope of validity.}
Applies iff \axopsurf{} = \texttt{multi-mode-with-OP\_SUCCESS}. The
only covenant with this axis value is CCV (BIP-443). CTV, OP\_VAULT,
and CAT+CSFS all have \axopsurf{} = \texttt{single-mode} and expose
a single semantic opcode with no mode selector.

\paragraph{(6) Rational-attacker check.}
Rational whenever a misconfigured vault exists on-chain. The
attacker's payoff is $V$ (full vault theft), the cost is a single
spend of $110\vB{} \times r$. The class is a developer-error surface,
not a cryptographic vulnerability.

\paragraph{(7) Counter-factual mitigations.}
Three distinct counter-factuals:
\begin{enumerate}
\item \axopsurf{} = \texttt{single-mode}. One opcode per primitive
  (as in the other three proposals) eliminates the class by
  construction.
\item Bitmask-gated mode check at script-parse time. Reject
  undefined modes before execution; eliminates the class without
  revising soft-fork semantics, but forfeits the forward-compatibility
  benefit.
\item Library-provided named constants. Makes the undefined-mode
  surface unreachable from application code; residual risk in
  independent implementations that hand-write the mode byte. This is
  the pymatt approach and is the recommended mitigation for CCV
  deployments.
\end{enumerate}

\paragraph{Framing note.}
The class is \emph{specified behaviour} under BIP-443's
forward-compatibility design (inherited from BIP-341/BIP-342), not a
newly-discovered vulnerability. The thesis contribution here is the
systematic mode sweep and the class-level framing that ties the
surface to axis \axopsurf{}, not the discovery of
\texttt{OP\_SUCCESSx} semantics.
